%% toc: P3 28 28.9
\section{Proofs for decidability by elimination and interpretation}
\label{sec:annex-p03-w03-p03-022-decidability}

\begin{annexitem}{annex:p03-w03-p03-022-one-step}{thm:p03-w03-p03-022-one-step}{An elimination step suffices}
\begin{proof}
Fix an elimination step \(\Sigma\) for \(T\). Define \(E(\varphi)\) by
recursion on \(\varphi\). For atomic \(\varphi\), and for \(\top\) and
\(\bot\), put \(E(\varphi)=\varphi\). Let \(E\) commute with each
propositional connective. Put \(E(\forall y\,\varphi)=E(\neg\exists y\,\neg\varphi)\).
For \(\exists y\,\varphi\), first compute the quantifier-free formula
\(\psi=E(\varphi)\). Compute a disjunctive normal form
\(\bigvee_{i<m}\theta_i\) of \(\psi\), in which each \(\theta_i\) is a
finite conjunction of literals of \(\psi\); this is a propositional
transformation of the atoms of \(\psi\), it introduces no variable, and
\(\psi\leftrightarrow\bigvee_i\theta_i\) is a tautology. Here \(m=0\) when
\(\psi\) is propositionally unsatisfiable, and the empty conjunction is
\(\top\). Put
\[
  E(\exists y\,\varphi)=\bigvee_{i<m}\psi_i,\qquad
  \psi_i=\begin{cases}\top,&\theta_i=\top,\\
  \Sigma(\exists y\,\theta_i),&\text{otherwise.}\end{cases}
\]
Each clause is a computable operation on codes, so \(E\) is computable,
and \(E(\varphi)\) is quantifier-free by construction.

We prove by induction on \(\varphi\) that
\(\operatorname{fv}(E(\varphi))\subseteq\operatorname{fv}(\varphi)\) and
\(T\entails\forall\bar x\,(\varphi\leftrightarrow E(\varphi))\). The atomic
case is trivial, and the propositional cases follow because truth of a
Boolean combination at an assignment depends only on the truth of its
components there. The universal case reduces to the negation and
existential cases. For \(\exists y\,\varphi\), let \(\mathcal M\models T\) and
let \(s\) be an assignment. By the induction hypothesis, for every
\(b\in M\), \(\mathcal M\models\varphi[s(y{:}b)]\) iff
\(\mathcal M\models\psi[s(y{:}b)]\), so \(\exists y\,\varphi\) and
\(\exists y\,\psi\) have the same truth value at \(s\). The equivalence
\(\exists y\bigvee_i\theta_i\leftrightarrow\bigvee_i\exists y\,\theta_i\) is
valid. Since every structure is nonempty, \(\exists y\,\top\) is valid, and
for the other \(i\) the step gives
\(T\entails\exists y\,\theta_i\leftrightarrow\psi_i\). Finally
\(\operatorname{fv}(\psi_i)\subseteq\operatorname{fv}(\exists y\,\theta_i)
\subseteq\operatorname{fv}(\psi)\setminus\{y\}
\subseteq\operatorname{fv}(\varphi)\setminus\{y\}
=\operatorname{fv}(\exists y\,\varphi)\).
\end{proof}
\end{annexitem}

\begin{annexitem}{annex:p03-w03-p03-022-decision}{thm:p03-w03-p03-022-decision}{Decision from elimination}
\begin{proof}
Let \(E\) be the elimination algorithm.

Item 1. Given a sentence \(\sigma\), compute \(\psi=E(\sigma)\). Since
\(\operatorname{fv}(\psi)\subseteq\operatorname{fv}(\sigma)=\emptyset\), \(\psi\)
is a quantifier-free sentence, and \(T\entails\sigma\leftrightarrow\psi\).
Therefore \(T\entails\sigma\) iff \(T\entails\psi\), and the latter is
decided by hypothesis.

Item 2. A quantifier-free sentence \(\psi\) is a Boolean combination of
atomic sentences and the constants \(\top,\bot\). By induction on \(\psi\),
either \(T\entails\psi\) or \(T\entails\neg\psi\). For atoms this is the
hypothesis, and \(\top,\bot\) are settled. If \(\psi=\neg\chi\), the two
alternatives for \(\chi\) exchange. If \(\psi=\chi_1\land\chi_2\), then
either both conjuncts are entailed, and so is \(\psi\), or some
\(\neg\chi_i\) is entailed, and then so is \(\neg\psi\); the other
connectives are definable from these. For an arbitrary sentence
\(\sigma\), \(T\entails\sigma\leftrightarrow E(\sigma)\), so \(T\) entails
\(\sigma\) or \(\neg\sigma\). If the signature has neither constant
symbols nor nullary relation symbols, every term contains a variable, so
every atomic formula has a free variable and the hypothesis on atomic
sentences is vacuous.

Item 3. Define \(v(\alpha)=1\) if \(T\entails\alpha\) and \(v(\alpha)=0\)
otherwise, for atomic sentences \(\alpha\); this is computable by
hypothesis. Extend \(v\) to quantifier-free sentences by truth tables. The
induction of item 2 shows that \(v(\psi)=1\) implies \(T\entails\psi\) and
\(v(\psi)=0\) implies \(T\entails\neg\psi\). As \(T\) is consistent, it
does not entail both \(\psi\) and \(\neg\psi\), so \(T\entails\psi\) iff
\(v(\psi)=1\). Thus the hypothesis of item 1 holds, and \(T\) is decidable;
it is complete by item 2. Let \(\mathcal A\models T\). Every consequence of
\(T\) is true in \(\mathcal A\). Conversely, if \(\mathcal A\models\sigma\),
then \(T\) does not entail \(\neg\sigma\), because \(\mathcal A\) is a model of \(T\) in
which \(\neg\sigma\) fails, so \(T\entails\sigma\) by completeness. Hence
\(\operatorname{Th}(\mathcal A)=\{\sigma:T\entails\sigma\}\).
\end{proof}
\end{annexitem}

\begin{annexitem}{annex:p03-w03-p03-022-dlo-step}{thm:p03-w03-p03-022-dlo-step}{Correctness and consequences for \(\mathsf{DLO}\)}
\begin{proof}
Let \(\mathcal M\models\mathsf{DLO}\) and fix values for the variables other
than \(y\). Each item replaces a formula by one with the same truth value in
\(\mathcal M\).

Item 1. In a linear order exactly one of \(u<v\), \(u=v\), \(v<u\) holds.
Hence \(\neg(u<v)\) is equivalent to \(v<u\lor u=v\), and \(\neg(u=v)\) to
\(u<v\lor v<u\). The existential quantifier distributes over the resulting
disjunction.

Item 2. \(y=y\) is valid, and \(y<y\) is unsatisfiable by irreflexivity,
so a conjunction containing it has no witness.

Item 3. For a variable \(x\) distinct from \(y\),
\(\exists y\,(y=x\land\chi(y))\leftrightarrow\chi(x)\) is valid.

Item 4. After items 1 to 3 and the removal of \(y\)-free conjuncts, the
only atoms containing \(y\) are \(x_i<y\) and \(y<z_j\) with \(x_i,z_j\)
distinct from \(y\). If \(b\) satisfies all of them, then
\(x_i<b<z_j\), and transitivity gives \(x_i<z_j\). Conversely suppose
\(x_i<z_j\) for all \(i\in L\), \(j\in U\). If both sets are nonempty, let
\(l\) be the largest of the values \(x_i\) and \(u\) the smallest of the
values \(z_j\); these exist in a linear order because the sets are finite
and nonempty, and \(l<u\) because \(l=x_i\) and \(u=z_j\) for some
\(i,j\). Density gives \(b\) with \(l<b<u\), and then
\(x_i\le l<b<u\le z_j\) for all \(i,j\). If \(U\) is empty, the absence of
a greatest element gives \(b>l\); if \(L\) is empty, the absence of a least
element gives \(b<u\); if both are empty, nonemptiness gives some \(b\).
In each case \(b\) witnesses the conjunction.

Every item is computable and introduces no free variable, so the method is
an elimination step, and \cref{thm:p03-w03-p03-022-one-step} gives
effective elimination. The signature \(\{<\}\) has no constant or nullary
relation symbols. The structure \((\mathbb Q;<)\) is a model: it is a
nonempty linear order, \(p<\tfrac{p+q}2<q\) whenever \(p<q\), and
\(p-1<p<p+1\). So \(\mathsf{DLO}\) is consistent, and items 2 and 3 of
\cref{thm:p03-w03-p03-022-decision}, whose atomic hypothesis is vacuous
here, give completeness, decidability and \(\operatorname{Th}(\mathbb Q;<)=
\{\sigma:\mathsf{DLO}\entails\sigma\}\).
\end{proof}
\end{annexitem}

\begin{annexitem}{annex:p03-w03-p03-022-succ}{thm:p03-w03-p03-022-succ}{Elimination and decision for successor}
\begin{proof}
We use three consequences of \(T_{\mathrm S}\).

(a) \(T_{\mathrm S}\entails S^ku=S^kv\leftrightarrow u=v\) for every \(k\).
This is induction on \(k\) using injectivity of \(S\).

(b) For \(k\ne l\), \(T_{\mathrm S}\entails\forall u\,S^ku\ne S^lu\). Say
\(k<l\). If \(S^ku=S^lu=S^k(S^{l-k}u)\), then (a) gives \(u=S^{l-k}u\),
which contradicts the axiom for \(n=l-k\ge1\). Similarly, for \(k<l\),
\(S^k0=S^l0\) gives \(0=S^{l-k}0=S(S^{l-k-1}0)\) by (a), contradicting
\(\forall x\,Sx\ne0\). So the elements \(S^k0\) are pairwise distinct and
every model of \(T_{\mathrm S}\) is infinite.

(c) For every \(K\),
\(T_{\mathrm S}\entails\forall z\,\bigl(\exists y\,S^Ky=z\leftrightarrow
\bigwedge_{j<K}z\ne S^j0\bigr)\). For \(K=0\) both sides are true. Assume
the claim for \(K\). Suppose \(z=S^{K+1}y\). Then \(z\ne0\) by the first
axiom. If \(z=S^j0\) with \(1\le j\le K\), then injectivity gives
\(S^Ky=S^{j-1}0\) with \(j-1<K\), contrary to the claim for \(K\).
Conversely suppose \(z\ne S^j0\) for all \(j\le K\). Since \(z\ne0\), the
third axiom gives \(w\) with \(Sw=z\). If \(w=S^j0\) for some \(j<K\), then
\(z=S^{j+1}0\) with \(j+1\le K\), which is excluded; so the claim for \(K\)
gives \(y\) with \(S^Ky=w\), and then \(S^{K+1}y=z\).

Now let \(\mathcal M\models T_{\mathrm S}\). Item 1 is justified by (a) and (b).
Item 2: for \(k\le K\), applying (a) with \(K-k\) gives
\(S^ky=S^lv\leftrightarrow S^Ky=S^{l+K-k}v\), and the same for the
negation. After item 2 every occurrence of \(y\) is in the subterm
\(S^Ky\), so the conjunction has the form \(\chi(S^Ky)\) where \(\chi(z)\)
contains \(z\) only bare. Item 3 is the equivalence
\[
 \exists y\,\chi(S^Ky)\leftrightarrow
 \exists z\,\bigl(\exists y\,S^Ky=z\land\chi(z)\bigr)\leftrightarrow
 \exists z\,\Bigl(\bigwedge_{j<K}z\ne S^j0\land\chi(z)\Bigr),
\]
the first step being valid and the second being (c). In item 4, every
\(z\)-literal is \(z=t\) or \(z\ne t\) with \(t\) free of \(z\), since its
other side is \(S^lv\) with \(v\) not \(y\). If some conjunct is \(z=t\), the
valid equivalence \(\exists z\,(z=t\land\eta)\leftrightarrow\eta(t)\)
applies. Otherwise the \(z\)-conjuncts say that \(z\) avoids finitely many
values, and since \(\mathcal M\) is infinite by (b) some element does so; the
existential is then true and equivalent to \(\top\). No item introduces a
free variable, and each is computable.

By \cref{thm:p03-w03-p03-022-one-step}, \(T_{\mathrm S}\) eliminates
quantifiers effectively. The atomic sentences are \(S^k0=S^l0\); by (b)
\(T_{\mathrm S}\) entails it when \(k=l\) and entails its negation when
\(k\ne l\), and which case holds is computable. The structure
\((\N;0,S)\) with \(S(n)=n+1\) satisfies the axioms: \(n+1\ne0\),
\(n+1=m+1\) implies \(n=m\), each \(n\ne0\) equals \((n-1)+1\), and
\(n+k\ne n\) for \(k\ge1\). Item 3 of \cref{thm:p03-w03-p03-022-decision}
therefore gives completeness, decidability and
\(\operatorname{Th}(\N;0,S)=\{\sigma:T_{\mathrm S}\entails\sigma\}\).

A sentence \(\sigma\) of \(\{S\}\) is also a sentence of \(\{0,S\}\), and
its truth in \((\N;0,S)\) does not depend on the interpretation of \(0\),
which it does not mention; so \((\N;S)\models\sigma\) iff
\(T_{\mathrm S}\entails\sigma\), which is decidable.
\end{proof}
\end{annexitem}

\begin{annexitem}{annex:p03-w03-p03-022-aba}{thm:p03-w03-p03-022-aba}{Elimination and decision for atomless Boolean algebras}
\begin{proof}
We use the following identities, valid in every Boolean algebra.
\begin{enumerate}
\item[(i)] \(u=v\) iff \(u\triangle v=0\). Indeed \(u\triangle v=0\) iff
  \(u\sqcap-v=0\) and \(-u\sqcap v=0\), that is, iff \(u\le v\) and \(v\le u\).
\item[(ii)] \(u\sqcup v=0\) iff \(u=0\) and \(v=0\); and \(u\le v\) iff
  \(u\sqcap-v=0\).
\item[(iii)] For every term \(\tau(\bar x,y)\),
  \(\tau=(\tau_1\sqcap y)\sqcup(\tau_0\sqcap-y)\), where
  \(\tau_1=\tau(\bar x,1)\) and \(\tau_0=\tau(\bar x,0)\).
\end{enumerate}
We prove (iii) by induction on \(\tau\). For \(\tau=y\) the right side is
\((1\sqcap y)\sqcup(0\sqcap-y)=y\). For \(\tau\) not containing \(y\) it is
\((\tau\sqcap y)\sqcup(\tau\sqcap-y)=\tau\). Write
\(\tau=(a\sqcap y)\sqcup(b\sqcap-y)\) and \(\tau'=(a'\sqcap y)\sqcup(b'\sqcap-y)\).
Distributivity and \(y\sqcap-y=0\) give
\(\tau\sqcap\tau'=(a\sqcap a'\sqcap y)\sqcup(b\sqcap b'\sqcap-y)\) and
\(\tau\sqcup\tau'=((a\sqcup a')\sqcap y)\sqcup((b\sqcup b')\sqcap-y)\).
For the complement, De Morgan and distributivity give
\(-\tau=(-a\sqcup-y)\sqcap(-b\sqcup y)
=(-a\sqcap-b)\sqcup(-a\sqcap y)\sqcup(-b\sqcap-y)\), and
\(-a\sqcap-b=(-a\sqcap-b\sqcap y)\sqcup(-a\sqcap-b\sqcap-y)\) lies below the
other two joinands, so \(-\tau=(-a\sqcap y)\sqcup(-b\sqcap-y)\). In each case
the coefficients are the values of the compound term at \(y=1\) and at
\(y=0\), which proves (iii).

Correctness of items 1 to 3 follows from (i) to (iii) and distributivity:
\(\tau=0\) iff \(\tau_1\sqcap y=0\) and \(\tau_0\sqcap-y=0\); a finite family
of equations \(\tau^{(i)}=0\) is equivalent to \(a\sqcap y=0\land b\sqcap-y=0\)
because \(\bigsqcup_i(\tau^{(i)}_1\sqcap y)=a\sqcap y\) and similarly for
\(b\); and \(\sigma\ne0\) iff \(\sigma_1\sqcap y\ne0\) or
\(\sigma_0\sqcap-y\ne0\). By (ii), \(a\sqcap y=0\) says \(y\le-a\) and
\(b\sqcap-y=0\) says \(b\le y\).

It remains to show, in a model of \(\mathsf{ABA}\), that
\[
 \exists y\Bigl(b\le y\le-a\land\bigwedge_{j\in J}c_j\sqcap y\ne0
   \land\bigwedge_{k\in K}d_k\sqcap-y\ne0\Bigr)
\]
holds iff \(a\sqcap b=0\), \(c_j\sqcap-a\ne0\) for all \(j\in J\), and
\(d_k\sqcap-b\ne0\) for all \(k\in K\).

Necessity. From \(b\le y\le-a\) we get \(b\le-a\), so \(a\sqcap b=0\).
Also \(c_j\sqcap y\le c_j\sqcap-a\) and, since \(-y\le-b\),
\(d_k\sqcap-y\le d_k\sqcap-b\); a nonzero element lies below each right side.

Sufficiency. Put \(u=-a\sqcap-b\). From \(a\sqcap b=0\) we get \(b\le-a\) and
\(a\le-b\), hence
\(-a=(-a\sqcap b)\sqcup(-a\sqcap-b)=b\sqcup u\) and likewise \(-b=a\sqcup u\).
Let \(J'\) be the set of \(j\) with \(c_j\sqcap b=0\) and \(K'\) the set of
\(k\) with \(d_k\sqcap a=0\). For \(j\in J'\), put \(e_j=c_j\sqcap u\); then
\(e_j=(c_j\sqcap b)\sqcup(c_j\sqcap u)=c_j\sqcap-a\ne0\). For \(k\in K'\), put
\(f_k=d_k\sqcap u\); similarly \(f_k=d_k\sqcap-b\ne0\).

Let \(G\) be the finite set of the \(e_j\) and \(f_k\), all below \(u\). For a
sign map \(\varepsilon:G\to\{+,-\}\), put
\(p_\varepsilon=u\sqcap\bigsqcap_{g\in G}g^{\varepsilon(g)}\), with \(g^+=g\)
and \(g^-=-g\). Distinct sign maps differ at some \(g\), so
\(p_\varepsilon\sqcap p_{\varepsilon'}\le g\sqcap-g=0\). Expanding
\(u=u\sqcap\bigsqcap_{g\in G}(g\sqcup-g)\) by distributivity gives
\(u=\bigsqcup_\varepsilon p_\varepsilon\). For \(g\in G\), since \(g\le u\),
\(g=\bigsqcup_\varepsilon(g\sqcap p_\varepsilon)\), and \(g\sqcap p_\varepsilon\)
is \(p_\varepsilon\) when \(\varepsilon(g)=+\) and \(0\) otherwise; so
\(g=\bigsqcup\{p_\varepsilon:\varepsilon(g)=+\}\). Let \(P\) be the set of
\(\varepsilon\) with \(p_\varepsilon\ne0\). For each \(\varepsilon\in P\), the
atomless axiom gives \(q_\varepsilon\) with \(q_\varepsilon\le p_\varepsilon\),
\(q_\varepsilon\ne0\) and \(q_\varepsilon\ne p_\varepsilon\). Put
\(z=\bigsqcup_{\varepsilon\in P}q_\varepsilon\) (so \(z=0\) when \(P\) is
empty) and \(y=b\sqcup z\).

Then \(b\le y\), and \(y\le-a\) because \(b\le-a\) and \(z\le u\le-a\).

Let \(j\in J\). If \(c_j\sqcap b\ne0\), then \(c_j\sqcap y\ge c_j\sqcap b\) is
nonzero. Otherwise \(j\in J'\) and \(e_j\ne0\), so some \(\varepsilon\in P\)
has \(\varepsilon(e_j)=+\), whence \(q_\varepsilon\le p_\varepsilon\le e_j\le c_j\)
and \(q_\varepsilon\le z\le y\); thus \(0\ne q_\varepsilon\le c_j\sqcap y\).

Let \(k\in K\); note \(-y=-b\sqcap-z\). If \(d_k\sqcap a\ne0\), then
\(a\le-b\) and, since \(z\le-a\), \(a\le-z\); so
\(d_k\sqcap a\le d_k\sqcap-y\), which is nonzero. Otherwise \(k\in K'\),
\(f_k\ne0\), and some \(\varepsilon\in P\) has \(p_\varepsilon\le f_k\). Put
\(w=p_\varepsilon\sqcap-q_\varepsilon\). If \(w=0\) then
\(p_\varepsilon\le q_\varepsilon\), so \(q_\varepsilon=p_\varepsilon\), which is
excluded; hence \(w\ne0\). Moreover \(w\le f_k\le d_k\), \(w\le u\le-b\), and
\(w\sqcap z=\bigsqcup_{\varepsilon'\in P}(w\sqcap q_{\varepsilon'})=0\): for
\(\varepsilon'\ne\varepsilon\), \(q_{\varepsilon'}\le p_{\varepsilon'}\) is disjoint
from \(p_\varepsilon\ge w\), and \(w\sqcap q_\varepsilon=0\) by definition of
\(w\). So \(w\le-z\), and \(0\ne w\le d_k\sqcap-b\sqcap-z=d_k\sqcap-y\).

Thus item 4 is correct; it is computable and introduces no free variable,
so the method is an elimination step and
\cref{thm:p03-w03-p03-022-one-step} applies.

For the decision, every closed term has value \(0\) or \(1\) in every
Boolean algebra, and the value is computed in the two-element algebra
\(\{0,1\}\), which is a subalgebra of every model of \(\mathsf{BA}\). So
\(\mathsf{ABA}\) entails each atomic sentence \(t=s\) whose two values agree,
and entails its negation otherwise, because \(0\ne1\) is an axiom. For
consistency, let \(\mathcal R\) be the set of finite unions of intervals
\([p,q)\) with rational \(0\le p<q\le1\), together with \(\emptyset\), under
union, intersection and complement in \([0,1)\). The intersection of two
such intervals is empty or such an interval, and the complement of a finite
union of them is a finite union of them, since it is the union of the gaps
between their sorted endpoints; so \(\mathcal R\) is a Boolean algebra of sets
with \(\emptyset\ne[0,1)\). A nonzero element contains some \([p,q)\), and
\([p,\tfrac{p+q}2)\) is a nonzero element below it that misses the point
\(\tfrac{p+q}2\) of it; so \(\mathcal R\models\mathsf{ABA}\). Item 3 of
\cref{thm:p03-w03-p03-022-decision} gives completeness and decidability.
\end{proof}
\end{annexitem}

\begin{annexitem}{annex:p03-w03-p03-022-presburger}{thm:p03-w03-p03-022-presburger}{Elimination and decision with congruence predicates}
\begin{proof}
All equivalences below are truth in \(\mathcal Z\) at every assignment.
Replacing a term by its normal term preserves values, since the normal
term is computed by the ring laws of \(\mathbb Z\).

Item 1. For integers, \(s=t\) iff \(|t-s|<1\), which is the stated
conjunction; \(s<t\) iff \(0<t-s\); \(\neg(0<t)\) iff \(t\le0\) iff
\(0<1-t\). If \(n\ge2\) and \(n\nmid t\), let \(\rho\in\{1,\ldots,n-1\}\)
be the remainder of \(t\); then \(n\mid t+(n-\rho)\) with
\(n-\rho\in\{1,\ldots,n-1\}\). Conversely, if \(n\mid t+r\) with
\(1\le r\le n-1\), then \(t\equiv-r\not\equiv0\pmod n\). For \(n=1\),
\(\mathrm D_1(t)\) always holds, so its negation is \(\bot\). Literals in
which \(y\) has coefficient \(0\) are moved outside.

Item 2. Let \(m=\ell/|k|\), a positive integer. Multiplying by \(m\)
preserves positivity, so \(0<ky+t\) iff \(0<m(ky+t)\), whose \(y\)-term is
\(\pm\ell y\). For divisibility, \(n\mid s\) iff \(nm\mid ms\) iff
\(nm\mid-ms\); taking \(\pm m=\ell/k\) gives
\(\mathrm D_n(ky+t)\leftrightarrow\mathrm D_{nm}(\ell y+(\ell/k)t)\).

Item 3. \(\exists y\,\chi(\ell y)\leftrightarrow
\exists z\,(\mathrm D_\ell(z)\land\chi(z))\): put \(z=\ell y\), and
conversely take \(y=z/\ell\). The literal \(0<z+t\) is \(-t<z\), the literal
\(0<-z+t\) is \(z<t\), and \(\mathrm D_N(z+s)\) is of the third form;
\(\mathrm D_\ell(z)\) is the case \(c=0\). Thus \(H\) is nonempty.

Item 4. Each \(m_h\) divides \(M\), so \(\mathrm D_{m_h}(z+c_h)\) and
\(\mathrm D_{m_h}(z'+c_h)\) agree whenever \(z\equiv z'\pmod M\).

Suppose \(J\) is nonempty. If \(\eta(b_j+r)\) holds, \(z=b_j+r\) is a
witness. Conversely let \(\eta(z)\) hold, and choose \(j_0\) with
\(b_{j_0}\) largest among the values \(b_j\). Then \(z-b_{j_0}\ge1\). Let
\(r\in\{1,\ldots,M\}\) satisfy \(r\equiv z-b_{j_0}\pmod M\); then
\(r\le z-b_{j_0}\). Put \(z'=b_{j_0}+r\). We have \(z'\le z<a_i\) for all
\(i\), \(z'>b_{j_0}\ge b_j\) for all \(j\), and \(z'\equiv z\pmod M\), so
\(\eta(z')\) holds, which is the disjunct for \(j_0\) and \(r\).

Suppose \(J\) is empty. If \(\eta(z)\) holds, take
\(r\in\{1,\ldots,M\}\) with \(r\equiv z\pmod M\); periodicity gives
\(\bigwedge_h\mathrm D_{m_h}(r+c_h)\). Conversely, if that conjunction holds for
some \(r\), choose a natural number \(N\) so large that
\(z=r-NM<a_i\) for each of the finitely many \(i\); periodicity gives the
congruences at \(z\), so \(\eta(z)\) holds.

Each item is computable, and the output mentions only \(y\)-free terms, so
the method is an elimination step for \(\operatorname{Th}(\mathcal Z)\). By
\cref{thm:p03-w03-p03-022-one-step} the elimination is effective. A closed
term has a computable integer value, and atomic sentences
\(s<t\), \(s=t\), \(\mathrm D_n(t)\) are decided by computing with integers;
so the quantifier-free sentences true in \(\mathcal Z\) form a decidable set,
and item 1 of \cref{thm:p03-w03-p03-022-decision}, applied to the complete
theory \(\operatorname{Th}(\mathcal Z)\), shows that it is decidable.

For \(\mathcal N\), let \(\varphi^{\ge0}\) be the \(\Lang_{\mathbb Z}\)-formula
obtained from an \(\Lang_{\N}\)-formula \(\varphi\) by replacing each
\(\exists v\,\chi\) with \(\exists v\,(\neg v<0\land\chi)\) and each
\(\forall v\,\chi\) with \(\forall v\,(\neg v<0\to\chi)\). The set \(\N\) is closed
under \(0\), \(1\) and \(+\) in \(\mathbb Z\), and the symbols of
\(\Lang_{\N}\) have the same meaning on it, so \(\mathcal N\) is a substructure
of the \(\Lang_{\N}\)-reduct of \(\mathcal Z\). An induction on \(\varphi\)
shows \(\mathcal N\models\varphi(\bar a)\) iff
\(\mathcal Z\models\varphi^{\ge0}(\bar a)\) for all \(\bar a\) in \(\N\):
atomic formulas have the same term values in a substructure, connectives
commute, and relativized quantifiers range exactly over \(\N\). This is the
translation of \cref{def:p03-w03-p03-022-nat-in-int} written without its
redundant term-value clauses.

Eliminate quantifiers from \(\varphi^{\ge0}\) in \(\mathcal Z\) to get a
quantifier-free \(\psi\). Rewrite each atom of \(\psi\) as an
\(\Lang_{\N}\)-atom with the same truth value at all natural-number
arguments. Normalize the relevant difference to \(\sum_ik_ix_i+c\). The atom
\(0<\sum_ik_ix_i+c\) becomes
\[
 \sum_{k_i<0}|k_i|x_i+\max(-c,0)\;<\;\sum_{k_i>0}k_ix_i+\max(c,0),
\]
the corresponding equation is treated in the same way, and
\(\mathrm D_n(\sum_ik_ix_i+c)\) becomes \(\mathrm D_n(\sum_ik'_ix_i+c')\)
with \(k'_i,c'\in\{0,\ldots,n-1\}\) the remainders of \(k_i,c\) modulo
\(n\). Each rewriting is valid over \(\mathbb Z\), and the new atoms contain
only nonnegative coefficients, so they are \(\Lang_{\N}\)-atoms whose truth
in \(\mathcal N\) and \(\mathcal Z\) agree on \(\N\). The result \(\psi'\) is
computed from \(\varphi\), has no new free variable, and satisfies
\(\mathcal N\models\forall\bar x\,(\varphi\leftrightarrow\psi')\). Closed
\(\Lang_{\N}\)-terms have computable values, so
\(\operatorname{Th}(\mathcal N)\) is decidable by item 1 of
\cref{thm:p03-w03-p03-022-decision}. A sentence of \(\{0,1,+,<\}\) is an
\(\Lang_{\N}\)-sentence with the same truth value in \(\N\), which decides
\(\operatorname{Th}(\N;0,1,+,<)\).
\end{proof}
\end{annexitem}

\begin{annexitem}{annex:p03-w03-p03-022-presburger-boundary}{thm:p03-w03-p03-022-presburger-boundary}{The congruence predicates are needed}
\begin{proof}
By induction on terms, every term in the single variable \(x\) built from
\(0,1,+\) takes the value \(kx+c\) at \(x\in\N\), for fixed
\(k,c\in\N\). An atom is therefore \(kx+c<k'x+c'\) or \(kx+c=k'x+c'\). If
\(k=k'\), its truth value does not depend on \(x\). If \(k\ne k'\), then for
\(x>|c'-c|\) the number \((k'-k)x+(c'-c)\) is nonzero with the sign of
\(k'-k\), so the truth value is constant for all such \(x\). A
quantifier-free \(\psi(x)\) is a Boolean combination of finitely many atoms,
so its truth value is constant for all \(x\) beyond the largest of these
bounds. The even numbers contain arbitrarily large elements and omit
arbitrarily large elements, so no such \(\psi\) defines them. The formula
\(\exists y\,(y+y=x)\) does define them in \(\mathcal N\), so
\(\operatorname{Th}(\N;0,1,+,<)\) does not eliminate quantifiers in
\(\{0,1,+,<\}\). Decidability is part of
\cref{thm:p03-w03-p03-022-presburger}.
\end{proof}
\end{annexitem}

\begin{annexitem}{annex:p03-w03-p03-022-acf-step}{thm:p03-w03-p03-022-acf-step}{Correctness of the step}
\begin{proof}
Let \(K\models\mathsf{ACF}\). Polynomials below lie in \(K[y]\), and
degrees add under multiplication there because \(K\) has no zero divisors.

\emph{Algebraic facts.} (F1) For \(f\in K[y]\) and \(\alpha\in K\) there is
\(h\in K[y]\) with \(f=(y-\alpha)h+f(\alpha)\), and \(\deg h=\deg f-1\) when
\(\deg f\ge1\). By induction on \(\deg f\): if \(\deg f\le0\) take \(h=0\);
if \(f\) has degree \(d\ge1\) and leading coefficient \(e\), then
\(f-ey^{d-1}(y-\alpha)\) has degree below \(d\), and adding
\(ey^{d-1}\) to the quotient it receives by induction gives a quotient for
\(f\) of degree \(d-1\) and a constant remainder; substituting
\(y=\alpha\) shows that the remainder is \(f(\alpha)\).
(F2) If \(\alpha_1,\ldots,\alpha_e\) are distinct roots of \(g\), then
\(g=(y-\alpha_1)\cdots(y-\alpha_e)w\) for some \(w\). By induction on
\(e\): by (F1), \(g=(y-\alpha_e)g_1\), and for \(r<e\) we have
\[
0=g(\alpha_r)=(\alpha_r-\alpha_e)\,g_1(\alpha_r),\qquad
\alpha_r-\alpha_e\ne0,
\]
so \(\alpha_1,\ldots,\alpha_{e-1}\) are roots of
\(g_1\). In particular a nonzero \(g\) has at most \(\deg g\) roots.
(F3) Every \(f\) of degree \(d\ge1\) with leading coefficient \(b\) equals
\(b(y-\alpha_1)\cdots(y-\alpha_d)\) for some \(\alpha_r\in K\). By induction
on \(d\): \(b^{-1}f\) is monic of degree \(d\), so it has a root \(\alpha\)
by the root axiom for \(d\); by (F1), \(f=(y-\alpha)h\) with \(\deg h=d-1\)
and leading coefficient \(b\), and \(h=b\) when \(d=1\).
(F4) \(K\) is infinite, by item 5 of \cref{thm:p03-w03-p03-020-field-algebra}.

\emph{Evaluation.} Evaluation at \(\bar a\) is a ring homomorphism
\(\mathbb Z[\bar x][y]\to K[y]\). So products, truncations and identities
between parametric polynomials pass to \(K[y]\); padding does not change
the evaluated polynomial; and the atom \(c=0\) holds at \(\bar a\) exactly
when \(c(\bar a)=0\), because a term denotes the value of its integer
polynomial.

\emph{Induction.} Write
\[
Z(F_1,\ldots,F_k;G)=\{y\in K:F_1(\bar a)(y)=\cdots=F_k(\bar a)(y)=0,\
G(\bar a)(y)\ne0\}.
\]
We prove that \(\operatorname{Elim}(F_1,\ldots,F_k;G)\) holds in \(K\) at
\(\bar a\) iff \(Z\ne\emptyset\), by induction on the pair \(\mu\) (sum of formal degrees
of the \(F_i\), \(k\)) in the lexicographic order of \(\N\times\N\), which is a
well-order. Each recursive call lowers \(\mu\): in item 2 the sum is
unchanged and \(k\) drops; a truncation lowers the sum by one; in item 4,
each \(F_j\) with \(j\ne i\) has formal degree \(p_j\ge q\) and is replaced by
one of formal degree \(q-1\), and there is at least one such \(j\).

Item 1. Here \(Z=\{y:G(\bar a)(y)\ne0\}\). If some coefficient \(c\) of
\(G\) has \(c(\bar a)\ne0\), then \(G(\bar a)\) is a nonzero polynomial with
finitely many roots by (F2), and \(K\) is infinite by (F4), so
\(Z\ne\emptyset\) and the output holds. If every \(c(\bar a)=0\), then
\(G(\bar a)\) is the zero polynomial, \(Z=\emptyset\) and the output fails.

Item 2. \(F_i(\bar a)\) is the constant \(c(\bar a)\). If
\(c(\bar a)=0\), dropping \(F_i\) does not change \(Z\), and the induction
hypothesis applies to the shorter list. If \(c(\bar a)\ne0\), then
\(Z=\emptyset\) and the output fails at its first conjunct.

Item 3. If \(b(\bar a)=0\), then \(F(\bar a)=F_{\le q-1}(\bar a)\), so \(Z\)
is unchanged by the truncation, the second disjunct fails, and the
induction hypothesis gives the claim. Let \(b(\bar a)\ne0\), and put
\(f=F(\bar a)\), of degree exactly \(q\ge1\), and \(g=G(\bar a)\). We show
\(Z\ne\emptyset\iff f\nmid g^q\iff R(\bar a)\ne0\); the last condition is
the truth of the second disjunct at \(\bar a\).

If \(f\mid g^q\) and \(f(y)=0\), then \(g(y)^q=0\), so \(g(y)=0\); hence
\(Z=\emptyset\). Conversely let \(Z=\emptyset\), so every root of \(f\) is a
root of \(g\). By (F3), \(f=b(\bar a)\prod_{r\le q}(y-\alpha_r)\). Let
\(\beta_1,\ldots,\beta_e\) be the distinct values among the \(\alpha_r\),
with multiplicities \(m_1,\ldots,m_e\); then
\(f=b(\bar a)\prod_s(y-\beta_s)^{m_s}\) and \(m_s\le q\). Each \(\beta_s\) is a
root of \(g\), so \(g=\prod_s(y-\beta_s)\,w\) by (F2), and
\[
g^q=\prod_s(y-\beta_s)^q\,w^q
=f\cdot b(\bar a)^{-1}\prod_s(y-\beta_s)^{q-m_s}\,w^q ,
\]
so \(f\mid g^q\). Thus \(Z\ne\emptyset\) iff \(f\nmid g^q\).

Let \(N=\max(qg_0,q-1)\), with \(g_0\) the formal degree of \(G\). Evaluating
\(b^{N-q+1}H=AF+R\) gives \(b(\bar a)^{N-q+1}g^q=A(\bar a)f+R(\bar a)\). If
\(f\mid g^q\), then \(f\mid R(\bar a)\); every nonzero multiple of \(f\) has
degree at least \(q\), while \(R(\bar a)\) has degree at most \(q-1\), so
\(R(\bar a)=0\). If \(R(\bar a)=0\), then \(f\) divides
\(b(\bar a)^{N-q+1}g^q\), and multiplying the quotient by
\(b(\bar a)^{-(N-q+1)}\) shows \(f\mid g^q\).

Item 4. If \(b(\bar a)=0\), the truncation does not change \(F_i(\bar a)\)
and the induction hypothesis applies. Let \(b(\bar a)\ne0\). For \(j\ne i\)
let \(p_j\) be the formal degree of \(F_j\); evaluating
\(b^{p_j-q+1}F_j=A_jF_i+F'_j\) at \(\bar a\) and at a point \(y\) with
\(F_i(\bar a)(y)=0\) gives
\(b(\bar a)^{p_j-q+1}F_j(\bar a)(y)=F'_j(\bar a)(y)\), so
\(F_j(\bar a)(y)=0\) iff \(F'_j(\bar a)(y)=0\). Hence
\(Z(F_1,\ldots,F_k;G)=Z(F'_1,\ldots,F'_k;G)\), the first disjunct fails,
and the induction hypothesis applies to the primed list.

\emph{The step.} Let \(\exists y\,\theta\) be given and prepared as in
\cref{def:p03-w03-p03-022-acf-step}. For \(\bar a\in K^m\) and \(y\in K\),
a literal \(t=s\) holds at \((\bar a,y)\) iff \((P_t-P_s)(\bar a,y)=0\),
and evaluating the coefficients first and then \(y\) gives the same value.
So \(\theta\) holds iff every \(F_i(\bar a)(y)=0\) and every
\(G_j(\bar a)(y)\ne0\), and the latter holds iff
\(G(\bar a)(y)=\prod_jG_j(\bar a)(y)\ne0\), because \(K\) has no zero
divisors and the empty product is \(1\ne0\). Hence
\(K\models\exists y\,\theta\) at \(\bar a\) iff \(Z\ne\emptyset\) iff the
output holds at \(\bar a\). The atoms of the output are \(c=0\) with
\(c\in\mathbb Z[\bar x]\), written with the variables of \(c\) only, so
the output has no free variable outside
\(\operatorname{fv}(\exists y\,\theta)\). Every operation used (the
polynomials \(P_t\), integer polynomial arithmetic, truncation, padding, the recursion of
\cref{thm:p03-w03-p03-037-prem}, and the passage from a polynomial to a
term) is computable, and the recursion stops by the decrease of \(\mu\). This
holds in every model of \(\mathsf{ACF}\), and every model of
\(\mathsf{ACF}_p\) is one, so the method is an elimination step for both
theories.
\end{proof}
\end{annexitem}

\begin{annexitem}{annex:p03-w03-p03-022-acf}{thm:p03-w03-p03-022-acf}{Decision for algebraically closed fields}
\begin{proof}
By \cref{thm:p03-w03-p03-022-acf-step,thm:p03-w03-p03-022-one-step},
\(\mathsf{ACF}\) and every \(\mathsf{ACF}_p\) eliminate quantifiers
effectively.

\emph{Characteristic.} In a field \(K\), the set
\(\{n\in\mathbb Z:n\cdot1=0\}\) is an ideal of \(\mathbb Z\). If it is not
\(\{0\}\), its least positive element \(m\) generates it, \(m\ne1\) because
\(1\ne0\), and \(m\) is prime: if \(m=uv\) with \(1<u,v<m\), then
\((u\cdot1)(v\cdot1)=0\) forces \(u\cdot1=0\) or \(v\cdot1=0\), against the
minimality of \(m\). If \(\chi_q\) holds for a prime \(q\), then \(m\mid q\),
so \(m=q\) and the ideal is \(q\mathbb Z\). If \(\neg\chi_q\) holds for every
prime \(q\), the ideal is \(\{0\}\). So in every model of \(\mathsf{ACF}_p\),
\(n\cdot1=0\) iff \(p\mid n\) when \(p\) is prime, and iff \(n=0\) when
\(p=0\); and every field satisfies \(\mathsf{Fld}_p\) for exactly one such
\(p\).

\emph{\(\mathsf{ACF}_p\).} An atomic sentence is \(t=s\) with \(t,s\) closed
terms, and \(P_t-P_s\) is an integer \(n\), computed from the terms; the
atom holds in a model iff \(n\cdot1=0\). By the previous paragraph,
\(\mathsf{ACF}_p\) entails the atom when \(p\mid n\) (resp.\ \(n=0\)) and
entails its negation otherwise, and which case holds is computable. The
field \(\mathbb Z/p\mathbb Z\) for \(p\) prime and the field \(\mathbb Q\)
for \(p=0\) have characteristic \(p\), and
\cref{thm:p03-w03-p03-020-closure-exists} embeds each in an algebraically
closed field of the same characteristic, a model of \(\mathsf{ACF}_p\).
So \(\mathsf{ACF}_p\) is consistent, and item 3 of
\cref{thm:p03-w03-p03-022-decision} shows that it is complete, decidable,
and axiomatizes \(\operatorname{Th}(K)\) for each of its models \(K\), which
are exactly the algebraically closed fields of characteristic \(p\).

\emph{\(\mathsf{ACF}\).} Each model of \(\mathsf{ACF}\) is a model of
\(\mathsf{ACF}_p\) for its characteristic \(p\), and
\(\mathsf{ACF}\subseteq\mathsf{ACF}_p\). Hence for every sentence
\(\psi\), \(\mathsf{ACF}\entails\psi\) iff \(\mathsf{ACF}_p\entails\psi\)
for every \(p\in\{0\}\cup\{\text{primes}\}\). For a quantifier-free
sentence \(\psi\), let \(v_p(\psi)\) be its truth value when each atom with
integer \(n\) is read as \(p\mid n\) (for \(p\) prime) or \(n=0\) (for
\(p=0\)); by the proof of item 3 of
\cref{thm:p03-w03-p03-022-decision}, \(\mathsf{ACF}_p\entails\psi\) iff
\(v_p(\psi)=1\). Let \(n_1,\ldots,n_r\) be the nonzero integers of the atoms
of \(\psi\) and \(N=\abs{n_1\cdots n_r}\) (\(N=1\) if \(r=0\)). For a prime
\(p\nmid N\), no \(n_i\) is divisible by \(p\), so \(v_p\) and \(v_0\) agree
on every atom and hence on \(\psi\). Therefore \(\mathsf{ACF}\entails\psi\)
iff \(v_0(\psi)=1\) and \(v_p(\psi)=1\) for the finitely many primes
\(p\mid N\), which are found by trial division. Given a sentence
\(\sigma\), compute its quantifier-free equivalent \(\psi\) and apply this
test: \(\mathsf{ACF}\) is decidable. It is not complete: \(\chi_2\) holds
in every model of \(\mathsf{ACF}_2\) and fails in every model of
\(\mathsf{ACF}_3\), and both theories are consistent.
\end{proof}
\end{annexitem}

\begin{annexitem}{annex:p03-w03-p03-022-rcf}{thm:p03-w03-p03-022-rcf}{Effective elimination and decision for real closed fields}
\begin{proof}
\emph{Computability of the construction.} Restrict the construction of
\cref{thm:p03-w03-p03-037-qe} to an input \(\exists x\,\theta\) with
\(\theta\) a conjunction of \(\Lang_{\mathrm{or}}\)-literals. Its stages are
the following, and each is computable.
\begin{enumerate}
\item Atoms. Each term \(t\) yields its integer polynomial \(P_t\) by
  recursion on \(t\); \(t_1=t_2\) becomes \(P_{t_2}-P_{t_1}=0\) and
  \(t_1<t_2\) becomes \(P_{t_2}-P_{t_1}>0\); \(\neg(P=0)\) and
  \(\neg(P>0)\) are expanded by trichotomy, and the result is put in
  disjunctive normal form. Each disjunct is a conjunction
  \(\bigwedge_{i\le s}\operatorname{sgn}P_i=\sigma_i\), and \(\exists x\) is
  distributed over the disjunction.
\item Sign tables. For one disjunct, \(\mathcal P=(P_1,\ldots,P_s)\) with
  the \(x\)-degrees as formal degrees and \(D\) their sum; the set
  \(\mathbb T_{s,D}\) of words of length at most \(2D+1\) over
  \(\{-1,0,1\}^s\) is finite and is listed. The output for the disjunct is
  the disjunction of the \(\theta_{\mathcal P,t}\) over the listed \(t\)
  having \(\sigma\) as a letter (\(0\ne0\) if there is none).
\item The formulas \(\theta_{\mathcal P,t}\). They are computed by
  recursion on \(\mu(\mathcal P)\) as in the annex proof of
  \cref{thm:p03-w03-p03-037-uniform-sign}. In the base case \(\theta\) is an
  explicit conjunction of sign conditions on coefficients, or \(0\ne0\).
  In the step, if the last member does not have the largest formal degree,
  move the first member of largest formal degree to the end; for the
  resulting bijection \(\iota\) and reordered list \(\mathcal P^\iota\),
  \cref{thm:p03-w03-p03-037-sublist} gives
  \(\operatorname{SIGN}(\mathcal P)=\rho_{\iota^{-1}}(\operatorname{SIGN}(\mathcal P^\iota))\),
  so \(\theta_{\mathcal P,t}\) is the disjunction of the
  \(\theta_{\mathcal P^\iota,u}\) over the listed \(u\) with
  \(\rho_{\iota^{-1}}(u)=t\). For a list whose last member has the largest
  formal degree, \(\mathcal P^\flat\), \(\mathcal R\) and \(\mathcal Q\) are
  formed by truncation, the formal derivative, extraction of coefficients
  and the pseudo-remainder recursion of \cref{thm:p03-w03-p03-037-prem},
  all of which are integer polynomial arithmetic, and
  \(\theta_{\mathcal P,t}\) is the displayed combination of
  \(\theta_{\mathcal P^\flat,t}\) and the \(\theta_{\mathcal Q,u}\) with
  \(u\) in the listed finite set \(\mathbb T_{\mathcal Q}\) and
  \(\Phi(u)=t\).
\item The maps \(\rho_\iota\) and \(\Phi\). \(\rho_\iota\) is given by an
  explicit rule on letters and positions. \(\Phi(u)\) is produced by steps
  (a)--(e) of the proof of the claim there, each of which is a finite
  search in \(u\) (the first member vanishing at a root, the largest index
  of a nonzero coefficient), a table look-up of a coordinate, an
  application of \(\rho\), a sign computation, or a choice among the four
  cases of \cref{thm:p03-w03-p03-037-interval} by the signs
  \(\tau_k,\tau_{k+1}\); and whether some step is undefined, in which case
  \(\Phi(u)=u\), is decided by the same finite inspection of \(u\).
\end{enumerate}
By induction on \(\mu(\mathcal P)\), using that each recursive call is on
\(\mathcal P^\flat\) or \(\mathcal Q\) with smaller \(\mu\) and that the
reordering is applied at most once before the step, the computation of all
\(\theta_{\mathcal P,t}\) halts. The atoms produced are sign conditions on
polynomials in \(\mathbb Z[\bar Y]\), where \(\bar Y\) lists the variables of
\(\theta\) other than \(x\), so the output has no free variable outside
\(\operatorname{fv}(\exists x\,\theta)\); and
\(\mathsf{RCF}\entails\exists x\,\theta\leftrightarrow\psi\) by the proof of
\cref{thm:p03-w03-p03-037-qe}. This is an elimination step for
\(\mathsf{RCF}\), so \cref{thm:p03-w03-p03-022-one-step} gives an
algorithm for all formulas; its recursion through connectives and
quantifiers is the reduction used in the proof of
\cref{thm:p03-w03-p03-037-qe}.

\emph{Atomic sentences.} In an ordered field, \(1\ne0\), and \(1<0\) would
give \(0<-1\) by adding \(-1\), then \(0<(-1)(-1)=1\), contradicting
\(1<0\); so \(0<1\). Hence \(k\cdot1>0\) for every integer \(k>0\), as a sum
of positive elements. An atomic sentence compares closed terms \(t,s\), whose
integer polynomials are integers \(n,m\) computed from the terms; in every
ordered field \(s-t\) has the value \((m-n)\cdot1\), which is positive,
zero or negative according as \(m-n\) is. So \(\mathsf{RCF}\) entails
\(t=s\) if \(n=m\) and its negation otherwise, and entails \(t<s\) if
\(n<m\) and its negation otherwise, and which case holds is computable.

\emph{Consequences.} \(\mathsf{RCF}\) has the model \(\mathbb R\), as in the
proof of \cref{thm:p03-w03-p03-037-complete}, so it is consistent. By item 3
of \cref{thm:p03-w03-p03-022-decision}, \(\mathsf{RCF}\) is complete and
decidable and axiomatizes \(\operatorname{Th}(K)\) for each real closed
field \(K\); with \(K=\mathbb R\) this decides
\(\operatorname{Th}(\mathbb R;0,1,+,-,\cdot,<)\). A sentence of
\(\{0,1,+,-,\cdot\}\) is an \(\Lang_{\mathrm{or}}\)-sentence with the same
truth value in \(\mathbb R\) and in its ring reduct, and membership in the
set of such sentences is decidable; so
\(\operatorname{Th}(\mathbb R;0,1,+,-,\cdot)\) is decided by the same
procedure.
\end{proof}
\end{annexitem}

\begin{annexitem}{annex:p03-w03-p03-022-transfer}{thm:p03-w03-p03-022-transfer}{Transfer of decidability and undecidability}
\begin{proof}
The translation of \cref{def:p03-imaginaries-translation} is defined by
recursion on syntax, using at each atom the domain, equality, relation and
graph formulas assigned to the symbols occurring there. When that
assignment is computable, so is \(\sigma\mapsto\sigma^I\); bound variables
are renamed by a fixed computable rule.

Since \(I\) is an interpretation with source theory \(\operatorname{Th}(\mathcal B)\)
and \(\mathcal B\models\operatorname{Th}(\mathcal B)\),
\cref{thm:p03-imaginaries-quotient-truth} applies to \(M=\mathcal B\). For a
sentence \(\sigma\) the assignment is empty, and the theorem gives
\(I(\mathcal B)\models\sigma\) iff \(\mathcal B\models\sigma^I\). An
isomorphism \(g:I(\mathcal B)\to\mathcal A\) preserves the truth of every
formula at corresponding assignments; this is an induction on formulas in
which the atomic case is the definition of isomorphism and the quantifier
case uses that \(g\) is a bijection. Hence \(\mathcal A\models\sigma\) iff
\(\mathcal B\models\sigma^I\).

The function \(\sigma\mapsto\sigma^I\) is therefore a computable many-one
reduction of \(\operatorname{Th}(\mathcal A)\) to \(\operatorname{Th}(\mathcal B)\).
If \(\operatorname{Th}(\mathcal B)\) is decidable, deciding
\(\sigma^I\in\operatorname{Th}(\mathcal B)\) decides
\(\sigma\in\operatorname{Th}(\mathcal A)\); the undecidability statement is the
contrapositive.
\end{proof}
\end{annexitem}
